\documentclass[11pt]{article}

% ---------------------------------------------------------------------------
% TERRA manuscript, rewritten against Andrew Modzelewski's review comments.
% See AJM_COMMENTS.md for the full comment list and how each is addressed.
%
% Editorial rules enforced in this draft:
%   * No LaTeX em-dashes ("---") anywhere in the prose  (C15).
%   * Every number traces to a real run; nothing is invented (lab policy).
%   * Prototype family is L1Md_T throughout               (C33, C158).
% ---------------------------------------------------------------------------

\usepackage[T1]{fontenc}
\usepackage{lmodern}
\usepackage[utf8]{inputenc}
\usepackage[margin=1in]{geometry}
\usepackage{graphicx}
\usepackage{booktabs}
\usepackage{amsmath}
\usepackage{amssymb}
\usepackage{xcolor}
\usepackage[numbers,sort&compress]{natbib}
\usepackage{caption}
\usepackage[colorlinks=true,allcolors=blue!55!black]{hyperref}

% Real schematics live in figs/. Dropping a populated reports8/ into this
% folder makes the data figures resolve with no edits to this file.
\graphicspath{{./}{figs/}{reports8/}{reports8/stage11_l1mdt/}}

% ---------------------------------------------------------------------------
% Figure slot: use the real file if present, otherwise draw a labelled box that
% names the missing file. Figures are never fabricated.
% ---------------------------------------------------------------------------
% NOTE: \IfFileExists does NOT consult \graphicspath, so every candidate
% directory must be probed explicitly here. Drop generated figures into any of
% figs/, reports8/ or reports8/stage11_l1mdt/ and they resolve automatically.
% Size figures by width and let the aspect ratio set the height, as a normal
% paper does. The previous version also capped height at 0.20\textheight, which
% fought the \resizebox wrapper at each call site: wide multi-panel figures hit
% the height cap, came out narrower than \linewidth, and were then scaled back
% up, so the effective size was arbitrary.
\newcommand{\fbinc}[2]{%
  \includegraphics[width=#1,keepaspectratio]{#2}}
\newcommand{\figorbox}[2][\linewidth]{%
  \IfFileExists{#2.pdf}{\fbinc{#1}{#2.pdf}}{%
  \IfFileExists{figs/#2.pdf}{\fbinc{#1}{figs/#2.pdf}}{%
  \IfFileExists{reports8/#2.pdf}{\fbinc{#1}{reports8/#2.pdf}}{%
  \IfFileExists{reports8/stage11_l1mdt/#2.pdf}{\fbinc{#1}{reports8/stage11_l1mdt/#2.pdf}}{%
  \IfFileExists{#2.png}{\fbinc{#1}{#2.png}}{%
  \IfFileExists{figs/#2.png}{\fbinc{#1}{figs/#2.png}}{%
  \IfFileExists{reports8/#2.png}{\fbinc{#1}{reports8/#2.png}}{%
  \IfFileExists{reports8/stage11_l1mdt/#2.png}{\fbinc{#1}{reports8/stage11_l1mdt/#2.png}}{%
    \setlength{\fboxsep}{0pt}%
    \fbox{\parbox[c][0.18\textheight][c]{\dimexpr#1-2\fboxrule\relax}{%
      \centering\ttfamily\detokenize{#2.pdf}\\[8pt]%
      \normalfont\itshape\small figure not yet generated; see README.md%
    }}%
  }}}}}}}}%
}

\newcommand{\cYes}{\textcolor{green!45!black}{\checkmark}}
\newcommand{\cPart}{\textcolor{orange!85!black}{$\circ$}}
\newcommand{\cNo}{\textcolor{black!30}{$\times$}}
\newcommand{\todo}[1]{\textcolor{red!70!black}{[\textbf{TODO:} #1]}}

% ---------------------------------------------------------------------------
% MEASURED VALUES FOR THE L1Md_T PROTOTYPE  (read this before editing numbers)
%
% How it works, in one rule: a generated file wins if it exists, otherwise the
% fallback below is used. Generated files use \providecommand, so \inputL runs
% FIRST and any value it defines takes precedence over the \providecommand
% fallbacks that follow.
%
% \inputL only ever looks in reports8/stage11_l1mdt/, i.e. the L1Md_T-specific
% output directory. This is deliberate: the *_measured_values.tex files sitting
% in the REPO ROOT are left over from an older IAP run (they say
% \grnaOTFamily{IAP}, 12577 copies, 25 lineages) and would silently overwrite
% the L1Md_T numbers with the wrong family's results. Restricting the search
% path means they can never be picked up, so there is nothing to remember.
%
% Therefore: when a fresh L1Md_T run lands, just drop reports8/ next to this
% file and rebuild. The numbers update themselves. Do not paste numbers into
% the prose; the prose only ever references these macros.
% ---------------------------------------------------------------------------
% Do NOT input reports8/measured_values.tex: the aggregate file the cross-family
% script writes is brace-unbalanced (it drops the closing "}" on values that wrap
% a \texttt{...}, e.g. \lOneTPhyloMasterName), which breaks the build. The
% per-family files below are well-formed and are the authoritative source.
\newcommand{\inputL}[1]{%
  \IfFileExists{reports8/stage11_l1mdt/#1}{\input{reports8/stage11_l1mdt/#1}}{%
    \IfFileExists{figs/#1}{\input{figs/#1}}{}}}
\inputL{phylo_measured_values.tex}
\inputL{grna_offtarget_measured_values.tex}
\inputL{transduction_measured_values.tex}
\inputL{antisense_measured_values.tex}
\inputL{cluster_validation_values.tex}
\inputL{cluster_search_values.tex}
\inputL{cluster_crosscheck_values.tex}

% Fallbacks matching the 2026-07-14 run, superseded automatically by the files above.
\providecommand{\phyloNCopies}{23639}
\providecommand{\phyloNIntact}{9388}
\providecommand{\phyloMasterName}{\texttt{chrX:40304675}}
\providecommand{\phyloMasterDiv}{5.35}
\providecommand{\phyloMedianAgeMyr}{9.02}
\providecommand{\grnaOTNGuides}{400}
\providecommand{\grnaOTNFrontier}{17}
\providecommand{\grnaOTBestGuide}{\texttt{AGGGAGGGAAGGAGGGAGGG}}
\providecommand{\transNLineages}{3}
\providecommand{\transNPolyaSignal}{8227}
\providecommand{\antiNBidir}{13246}

% Fallbacks: the real L1Md_T values from the run Andrew reviewed. Each is
% superseded automatically if the corresponding generated file above is present.
\newcommand{\protoFamily}{L1Md\_T}
\newcommand{\protoAssembly}{mm10}
\newcommand{\protoLoci}{23,639}
\newcommand{\protoClusters}{\csClusters}
\newcommand{\protoExprLoci}{7,490}
\newcommand{\protoMedianDiv}{3.5}
% 6 of 6 real developmental stages are significant. The earlier draft said 8:
% that run also tested expr_start/expr_stop, which are genomic coordinates rather
% than expression (fixed in run_line_sine_ltr_analysis.py _NON_EXPR), so two of
% the eight "stages" only reflected the clusters sitting at different positions.
\newcommand{\protoDESig}{\csDeSig}
\newcommand{\protoDETotal}{\csDeTotal}
\newcommand{\protoStages}{\csNStages}
\newcommand{\protoDivisor}{5}
% n_neighbors=15 per the run's cache_l1md_t_meta.json (the reviewed draft said 30,
% which does not match the run whose figures this paper now ships).
\newcommand{\protoNeighbors}{\csNeighbors}
\newcommand{\protoKmer}{\csKmer}
\newcommand{\protoMinClusterSize}{\csMinClusterSize}
\newcommand{\protoMinSamples}{\csMinSamples}
% Noise class: HDBSCAN leaves 12,720 of 23,639 loci unassigned (53.8%); the two
% clusters therefore hold 10,919 loci. Measured from cache_l1md_t_clustered.pkl.
\newcommand{\protoNoisePctMeasured}{\csNoisePct}
% the superseded default: min_cluster_size=N/5 gave 2 clusters and 53.8% noise
\newcommand{\oldClusters}{2}
\newcommand{\oldNoisePct}{53.8}
\newcommand{\oldMinClusterSize}{4,727}
\newcommand{\oldMedianARI}{0.21}
% copy-resolved module results (L1Md_T)
% Module numbers: alias the generated macros so a fresh run updates them itself.
% (The reviewed draft named chr1:95511396 as the master element; the run whose
% figures this paper ships nominates chrX:40304675, so the generated value wins.)
\newcommand{\protoIntact}{\phyloNIntact}
\newcommand{\protoMaster}{\phyloMasterName}
\newcommand{\protoMasterDiv}{\phyloMasterDiv}
\newcommand{\protoMedianAge}{\phyloMedianAgeMyr}
\newcommand{\protoGuides}{\grnaOTNGuides}
\newcommand{\protoFrontier}{\grnaOTNFrontier}
\newcommand{\protoBestGuide}{\grnaOTBestGuide}
\newcommand{\protoLineages}{\transNLineages}
\newcommand{\protoPolyA}{\transNPolyaSignal}
\newcommand{\protoBidir}{\antiNBidir}
% supplementary families
\newcommand{\bOneMTLoci}{70,685}
\newcommand{\bOneMTExprLoci}{21,407}
\newcommand{\bOneMTMedianDiv}{7.5}
\newcommand{\iapLoci}{1,485}
\newcommand{\iapExprLoci}{1,426}
\newcommand{\iapMedianDiv}{1.6}
% Clustering / noise validation (C2). Populated by run_cluster_validation.py;
% "(pending run)" until that script has been run for this family.
\providecommand{\xcAccuracy}{36.1}\providecommand{\xcChance}{20.0}
\providecommand{\xcDivOwn}{0.035}\providecommand{\xcDivOther}{0.030}
\providecommand{\xcSeparation}{-0.005}\providecommand{\xcNTested}{599}
\providecommand{\xcSample}{120}\providecommand{\xcWorstCluster}{7.5}
\newcommand{\intactRich}{91.2}\newcommand{\intactPoor}{0.1}
\providecommand{\csClusters}{5}\providecommand{\csNoisePct}{7.2}
\providecommand{\csKmer}{5}\providecommand{\csNeighbors}{50}
\providecommand{\csMinClusterSize}{2000}\providecommand{\csMinSamples}{25}
\providecommand{\csNConfigs}{288}\providecommand{\csNInTarget}{44}
\providecommand{\csDeSig}{54}\providecommand{\csDeTotal}{60}
\providecommand{\csNStages}{6}\providecommand{\csReproARI}{(pending run)}
\providecommand{\cvNLoci}{(pending run)}\providecommand{\cvNConfigs}{(pending run)}
\providecommand{\cvRefKmer}{6}\providecommand{\cvRefNeighbors}{15}
\providecommand{\cvNoisePct}{(pending run)}
\providecommand{\cvNoiseMinPct}{(pending run)}\providecommand{\cvNoiseMaxPct}{(pending run)}
\providecommand{\cvMinARI}{(pending run)}\providecommand{\cvMedianARI}{(pending run)}
\providecommand{\cvMinARIcore}{(pending run)}
\providecommand{\cvClustersMin}{(pending run)}\providecommand{\cvClustersMax}{(pending run)}

\captionsetup{font=small,labelfont=bf}
\setlength{\parskip}{2pt}
\setlength{\emergencystretch}{3em}

\title{\textbf{TERRA}: an accessible, HPC-integrated pipeline for downstream
transposable-element family analysis and CRISPR reagent design}
% Author list is not final (C1): add Vamshi's full name and affiliation, the
% graduate students, and anyone else who gave meaningful input, each with an
% affiliation. To be completed once the lab confirms the contributor list.
% Name (C0): "TERRA" (Transposable Element Repeat Resolved Analysis, formerly
% GAMECA) is used throughout; the final name is still to be crowdsourced with the
% lab, and a rename is a single find/replace.
\author{Anmol Dash \and Vamshidhar Nallamalla \and Katie Toohill \and Diego Jimenez \and Hyun An \and Andrew J. Modzelewski}
\date{\today}

\begin{document}
\maketitle

% ===========================================================================
\begin{abstract}
\noindent
Transposable elements (TEs) make up a large fraction of mammalian genomes and are
increasingly recognised as functional contributors to genome regulation in development and
disease~\citep{bourque2018}. Mature tools discover TE families from raw
assemblies~\citep{flynn2020} and quantify repeat
expression~\citep{yang2019squire,jin2015tetranscripts,bendall2019telescope}, but what happens
once a family is named is under-served: deciding which copies matter and turning that into
reagents a bench scientist can order. We present \textbf{TERRA}, a modular open-source
pipeline that takes a TE family name and a reference assembly and returns two reagent
deliverables, \textbf{family- and locus-specific primers} and \textbf{allele-aware CRISPRa/i
guide RNAs} scored against the full copy landscape on a coverage-versus-off-target frontier,
together with the evidence needed to interpret them. Because averaging across a family mixes a
minority of expressed, intact copies with a majority of inert relics, TERRA clusters copies
without alignment, separates the \emph{expressed} from the \emph{common} consensus, and layers
expression, motif, phylogenetic and structural evidence on each sub-group, emitting figures and
tables at every stage: the cluster embedding, per-cluster expression and differential-expression
tests, MAFFT alignments and consensus sequences, JASPAR motif enrichment, and a molecular-clock
phylogeny that nominates a master element. On the prototype family \emph{\protoFamily} in
\texttt{\protoAssembly}, \protoLoci{} annotated loci reduce to \protoClusters{} clusters,
\protoExprLoci{} carry SQuIRE expression~\citep{yang2019squire}, \protoIntact{} copies retain
ORF integrity, and the pipeline nominates a master-element candidate (\protoMaster) and
\protoGuides{} guides of which \protoFrontier{} are non-dominated. The same run executes
unchanged on a laptop or on LSF/Slurm clusters through Nextflow
DSL2~\citep{ditommaso2017,ewels2020}. TERRA is open source.
\end{abstract}

% ===========================================================================
\section{Introduction}

Interspersed repeats derived from TEs account for roughly half of the human
genome~\citep{lander2001} and a comparable or larger share of many other eukaryotic
genomes~\citep{bourque2018}, where they shape genome organisation, supply regulatory
sequence, and influence host gene expression in development and
disease~\citep{wicker2007,chuong2017enhancers}. The natural unit of analysis is the
\emph{family}, which groups copies descended from a common ancestral element.

The bioinformatics of TEs has historically concentrated on two upstream problems: building
libraries of family models \emph{de novo} from an assembly, and annotating a genome against
them. Tools such as RepeatModeler2~\citep{flynn2020} and the RepeatMasker/Dfam
ecosystem~\citep{hubley2016} address both well. A second, well-populated group quantifies
repeat expression, including SQuIRE~\citep{yang2019squire},
TEtranscripts/TElocal~\citep{jin2015tetranscripts}, Telescope~\citep{bendall2019telescope},
scTE~\citep{he2021scte} and ERVmap~\citep{tokuyama2018ervmap}.

The under-served workflow begins once a family has been named and quantified, and two problems
sit at its centre. First, which copies matter: a named family bundles a few expressed, intact
insertions with many truncated, inert relics, so any family-level analysis dilutes a signal
carried by a minority of loci. Resolving the family into sub-groups and asking which are
expressed is what makes the result interpretable, and it produces the evidence a biologist
reads directly (cluster embedding, per-cluster expression, alignments, consensus sequences,
motif enrichment). Second, the reagents: most TE projects end in an experiment, and
repetitiveness makes the necessary primers and guides hard to design, since a guide inevitably
matches many copies and must trade coverage of the family against off-target hits. TERRA treats
primers and guides as first-class outputs, reports the guide trade-off as a Pareto frontier
rather than one score, and returns both the reagents and the evidence behind them as figures and
tables rather than leaving the user to assemble them.

TERRA applies the philosophy of integrated analysis
platforms~\citep{galaxy2022,genepattern2006,ewels2020,ditommaso2017}, namely hide the
infrastructure, keep every step reproducible, and run anywhere. It is not a TE-discovery tool
and does not compete with RepeatModeler2~\citep{flynn2020}: it begins with a defined family and
assembly and automates the characterisation and reagent design that follow, with first-class
support for the LSF and Slurm schedulers that dominate academic HPC.

Figure~\ref{fig:journey} summarises the path this paper follows end to end on a single
prototype family.

\begin{figure}[htbp]
  \centering
  \figorbox[0.74\textwidth]{fig_prototype_journey}
  \caption{The TERRA path, from database annotation to ordered reagents, illustrated for the
    prototype family used throughout this paper. RepeatMasker/Dfam loci are sub-grouped
    without alignment, filtered by locus-level expression to separate copies that matter from
    inert relics, characterised by consensus and motif content, and converted into the two
    deliverables: family/locus-specific primers and allele-aware CRISPRa/i guides. This panel
    is a schematic of the workflow; every measured result it refers to is reported in
    Section~\ref{sec:results}.}
  \label{fig:journey}
\end{figure}

% ===========================================================================
\section{Results: \protoFamily{} from annotation to reagents}
\label{sec:results}

We use one family, \emph{\protoFamily}, as a prototype throughout: a young, active mouse
LINE-1 subfamily that is large, contains both intact and heavily truncated copies, and is
expressed in early development, so every stage is exercised. All results come from a single
\texttt{\protoAssembly} run with locus-level expression from SQuIRE~\citep{yang2019squire}
across six early-development stages (oocyte, pronucleus, 2-, 4-, 8-cell, morula). A SINE
(\emph{B1\_Mus2}) and an LTR (\emph{IAPLTR1\_Mm}) family, showing the workflow generalises
across TE classes, are in Section~\ref{sec:supp} and the Supplementary Files.

\subsection{Starting point: the annotated family}

TERRA takes the family name and assembly and retrieves loci and per-locus divergence from the
RepeatMasker track~\citep{hubley2016}, then extracts sequence from the reference. For
\emph{\protoFamily} in \texttt{\protoAssembly} this yields \protoLoci{} annotated loci. This
is the number a biologist starts with, and on its own it is not actionable: it says nothing
about which of those \protoLoci{} insertions are intact or worth targeting.

\subsection{Resolving copies rather than lumping them}
\label{sec:resolve}

Clustering is alignment-free. Each locus sequence is encoded as a $k$-mer count vector
(default $k{=}6$, giving a $4^6 = 4096$-dimensional sparse vector built with
scikit-learn~\citep{pedregosa2011}), reduced to 50 components by truncated SVD, embedded in
two dimensions with UMAP~\citep{mcinnes2018umap}, and partitioned with
HDBSCAN~\citep{mcinnes2017hdbscan}. This avoids a family-wide multiple alignment and scales
to tens of thousands of loci. For \emph{\protoFamily} the parameter search
(Section~\ref{sec:validation}) converged to a divisor of $N/\protoDivisor{}$ with UMAP
\texttt{n\_neighbors} $=\protoNeighbors{}$, $k=\protoKmer{}$, minimum cluster size
\protoMinClusterSize{} and \texttt{min\_samples} $=\protoMinSamples{}$, partitioning the family
into \protoClusters{} clusters with \protoNoisePctMeasured\% of loci left in HDBSCAN's noise
class.

These parameters were not taken on faith. The pipeline's heuristic minimum cluster size
$\lfloor N/\protoDivisor \rfloor$ demands \oldMinClusterSize{} loci per cluster for
\emph{\protoFamily}, which only \oldClusters{} groups can satisfy, forcing \oldNoisePct\% of the
family into the noise class; that is a property of the threshold, not the family. Searching
$k$, \texttt{n\_neighbors}, \texttt{min\_cluster\_size} and \texttt{min\_samples}
(\csNConfigs{} configurations, \texttt{run\_cluster\_search.py}), \csNInTarget{} recovered
\protoClusters{}--6 clusters, and the selected one cuts the unassigned fraction from
\oldNoisePct\% to \protoNoisePctMeasured\% (Figure~\ref{fig:cluster_search}). The remaining
noise class is kept as a labelled group but excluded from consensus building
(Section~\ref{sec:validation}), so unassigned loci cannot contaminate a consensus that reagents
target.

The payoff is shown in Figure~\ref{fig:filtering}: a family-level mean is dominated by the
inert majority, burying a real effect in the expressed, intact minority. Resolving copies and
intersecting them with expression separates these populations, and the consensus built from
expressed copies (the \emph{expressed} consensus) is what reagents should target, which is not
in general the consensus built from all annotated copies.

Of the \protoLoci{} loci, \protoExprLoci{} (roughly a third) carry SQuIRE expression entries,
summarised per cluster (units in Section~\ref{sec:validation}). Pairwise Wilcoxon rank-sum tests
with Benjamini--Hochberg correction find \protoDESig{} of \protoDETotal{} cluster-pair by stage
comparisons significant at $q<0.05$ across the \protoStages{} stages and \protoClusters{}
clusters. Expression is uneven: it concentrates in a minority of clusters and peaks at the
8-cell stage while the rest are near-silent (Figure~\ref{fig:proto_results}b), exactly the
separation a family-level average would remove. Figure~\ref{fig:proto_results} shows the
partition, the per-cluster expression profile and the divergence distribution.


\begin{figure}[htbp]
  \centering
  \figorbox[0.80\textwidth]{fig_line_results}
  \caption{\emph{\protoFamily} (\texttt{\protoAssembly}) analysed end to end in a single TERRA
    run. \textbf{(a)} UMAP embedding of $k$-mer-encoded loci coloured by HDBSCAN cluster.
    \textbf{(b)} Per-cluster mean SQuIRE expression (log1p) across the six developmental
    stages, computed from the \protoExprLoci{} loci carrying expression entries.
    \textbf{(c)} CpG-corrected Kimura divergence histogram, median \protoMedianDiv\%.}
  \label{fig:proto_results}
\end{figure}


\subsection{Age, intactness, and the master element}

Copies are aligned with MAFFT~\citep{katoh2002mafft}, a majority-rule consensus is built, and
each copy's Kimura two-parameter divergence~\citep{kimura1980} from that consensus is
converted to a molecular-clock age. ORF integrity separates intact from degraded copies, and a
master-element score combining intactness and divergence nominates the youngest, most intact
copies, which are the ones most likely to still be capable of mobilisation. For
\emph{\protoFamily}, \protoIntact{} of \protoLoci{} copies retain ORF integrity, the median
copy age is \protoMedianAge{} Myr, and the top master-element candidate is \protoMaster{} at
\protoMasterDiv\% divergence from the consensus (Figures~\ref{fig:phylo_tree}
and~\ref{fig:phylo_master}). 


\subsection{Deliverable 1: allele-aware CRISPRa/i guides}
\label{sec:crispr}

Designing CRISPR reagents against a repetitive family is the problem most tools leave to the
user. A guide has no single on-target site: it matches some number of family copies and some
elsewhere, so the design question is bi-objective (maximise on-family coverage, minimise
off-target hits) and cannot be reduced to one score without hiding the trade-off. TERRA
enumerates PAM-aware candidate guides from the consensus and per-copy sequence, scores each
against the \emph{full} copy landscape (and an optional background FASTA) with seed-anchored,
mismatch-tolerant matching, and reports the coverage-versus-off-target Pareto frontier. For
\emph{\protoFamily}, \protoGuides{} guides scored against all \protoLoci{} copies yield
\protoFrontier{} non-dominated design choices (Figure~\ref{fig:grna_pareto}), the best being
\protoBestGuide{}; selecting from the frontier is choosing a documented position on the
trade-off rather than trusting an opaque ranking. The same output supports CRISPRa and CRISPRi,
which depend on the guide-to-copy mapping rather than the effector
(Figure~\ref{fig:crispr}).


\begin{figure}[htbp]
  \centering
  \figorbox[0.56\textwidth]{fig_grna_offtarget_pareto}
  \caption{Measured coverage-versus-off-target frontier for \emph{\protoFamily}.
    \protoGuides{} candidate guides scored against \protoLoci{} copies; \protoFrontier{}
    guides are non-dominated. Each frontier point is a defensible design choice at a stated
    coverage and off-target count.}
  \label{fig:grna_pareto}
\end{figure}

\subsection{Deliverable 2: family- and locus-specific primers}

Primers are proposed from frequent $k$-mers in the (expressed) consensus and each candidate is
checked for specificity genome-wide across all chromosomes. Because the specificity check is
the most I/O-intensive step in the pipeline, it motivates the genome-cache engineering
described in Section~\ref{sec:design}. Primer candidates are emitted as a table with their
genome-wide match counts, and are cross-referenced against the per-locus expression so that a
user can select primers that report on the copies that are actually transcribed rather than on
the family as annotated.

\subsection{Additional copy-resolved evidence}

Three further modules were exercised on \emph{\protoFamily} and are reported here because they
bear directly on reagent design. 3' transduction detection groups copies sharing
downstream-tail $k$-mers absent from the consensus, alongside poly-A signal detection, and
resolves \protoLineages{} candidate lineages with \protoPolyA{} copies carrying an AATAAA
signal (Figure~\ref{fig:transduction}). Antisense/bidirectional promoter scanning of each 5'
region on both strands flags \protoBidir{} of \protoLoci{} copies as bidirectional
(Figure~\ref{fig:antisense}), which is relevant because a CRISPRa guide placed in a
bidirectional promoter may drive transcription in both directions. A further set of modules
(CTCF/TAD proximity, epigenetic overlay, orthologous-insertion calling, multi-assembly
liftover, ColabFold structure prediction, LTR structural classification, and automatic
subfamily resolution) place the same resolved copies in regulatory, evolutionary and
structural context, and are described in Section~\ref{sec:exploratory}.



% ===========================================================================
\section{How TERRA relates to existing tools}
\label{sec:comparison}

TERRA's neighbours fall into three groups (Table~\ref{tab:comparison}), and it consumes the
output of the first two rather than competing with them. \emph{Discovery and annotation:}
RepeatModeler2~\citep{flynn2020} and the RepeatMasker/Dfam ecosystem~\citep{hubley2016} produce
the locus coordinates TERRA starts from. \emph{Expression quantification:} SQuIRE~\citep{yang2019squire},
TEtranscripts/TElocal~\citep{jin2015tetranscripts}, Telescope~\citep{bendall2019telescope},
scTE~\citep{he2021scte}, ERVmap~\citep{tokuyama2018ervmap} and TE-Seq~\citep{kelsey2025teseq}
assign reads to repeats; TERRA consumes their per-locus counts (SQuIRE here) and does not
requantify. These answer ``how much is expressed''; TERRA addresses which copies to group,
what their consensus is, and what reagent targets them. \emph{General platforms:}
Galaxy~\citep{galaxy2022}, GenePattern~\citep{genepattern2006} and
nf-core/Nextflow~\citep{ewels2020,ditommaso2017} are portable substrates but ship no ready-made
TE-family workflow; TERRA builds on them, shipping a Nextflow DSL2 implementation importable as
a subworkflow. Table~\ref{tab:comparison} compares \emph{scope and access} from documented
capabilities; no cross-tool runtime is claimed.

\begin{table}[h]
  \centering
  \footnotesize
  \setlength{\tabcolsep}{5pt}
  \renewcommand{\arraystretch}{1.2}
  \begin{tabular}{@{}l cccccc@{}}
    \toprule
    \textbf{Tool} &
      \shortstack{TE-family\\focus} &
      \shortstack{Copy-resolved\\character.} &
      \shortstack{Reagent\\design} &
      \shortstack{Graphical\\interface} &
      \shortstack{Built-in\\LSF/Slurm} &
      \shortstack{Reprod.\\environ.} \\
    \midrule
    \textbf{TERRA} (this work)                  & \cYes & \cYes  & \cYes  & \cYes & \cYes  & \cYes  \\
    RepeatModeler2~\citep{flynn2020}            & \cYes & \cNo   & \cNo   & \cNo  & \cNo   & \cPart \\
    RepeatMasker/Dfam~\citep{hubley2016}        & \cYes & \cNo   & \cNo   & \cNo  & \cNo   & \cPart \\
    SQuIRE~\citep{yang2019squire}               & \cYes & \cPart & \cNo   & \cNo  & \cNo   & \cPart \\
    TEtranscripts/TElocal~\citep{jin2015tetranscripts} & \cYes & \cPart & \cNo & \cNo & \cNo & \cPart \\
    Telescope~\citep{bendall2019telescope}      & \cYes & \cPart & \cNo   & \cNo  & \cNo   & \cPart \\
    scTE~\citep{he2021scte}                     & \cYes & \cPart & \cNo   & \cNo  & \cNo   & \cPart \\
    ERVmap~\citep{tokuyama2018ervmap}           & \cYes & \cPart & \cNo   & \cNo  & \cNo   & \cPart \\
    TE-Seq~\citep{kelsey2025teseq}              & \cYes & \cPart & \cNo   & \cNo  & \cPart & \cPart \\
    Galaxy~\citep{galaxy2022}                   & \cNo  & \cPart & \cNo   & \cYes & \cPart & \cYes  \\
    GenePattern~\citep{genepattern2006}         & \cNo  & \cPart & \cNo   & \cYes & \cPart & \cYes  \\
    nf-core/Nextflow~\citep{ewels2020}          & \cNo  & \cPart & \cNo   & \cNo  & \cYes  & \cYes  \\
    \bottomrule
  \end{tabular}
  \caption{Capability comparison on documented features. \cYes~built-in; \cPart~partial, or
    possible but the user must assemble or configure it; \cNo~not provided.
    \emph{Copy-resolved characterisation}: partitioning a named family into sub-groups and
    building per-group consensus, as distinct from quantifying expression per locus.
    \emph{Reagent design}: primers and/or allele-aware CRISPR guides emitted for the family.
    Expression quantifiers are marked \cPart{} for copy-resolved characterisation because they
    report per-locus or per-family counts without grouping copies or building consensus. No
    cross-tool runtime is claimed.}
  \label{tab:comparison}
\end{table}

% ===========================================================================
\section{Validation and reporting conventions}
\label{sec:validation}

Expression is consumed as per-locus counts from an external quantifier
(SQuIRE~\citep{yang2019squire} throughout); TERRA does not requantify. Per-cluster values are
means of $\log_{1p}$ per-locus counts over the loci in that cluster carrying an entry; loci with
no entry are excluded rather than treated as zero, which is why \protoExprLoci{} of \protoLoci{}
loci enter the analysis.  Motif results are counts of JASPAR matches~\citep{castromondragon2022jaspar} per locus intersected with
BEDTools~\citep{quinlan2010bedtools}; enrichment is a per-cluster Fisher exact test on
loci-carrying-a-match, so a locus with several matches for one motif contributes once.

HDBSCAN labels loci it cannot confidently place as $-1$. These are
retained and reported as a labelled group but are excluded from per-cluster consensus building,
so a locus never contributes to a consensus it was not assigned to. The adopted configuration
leaves \protoNoisePctMeasured\% of loci unassigned.

The partition is a stratification of copy state, not a subfamily assignment; three measurements
bound what it can be claimed to be, reported with their negative results
(Appendix~\ref{sec:appendix}). \emph{Parameter dependence:} across \cvNConfigs{} configurations
of $k$ and \texttt{n\_neighbors} the cluster count ranged \cvClustersMin{}--\cvClustersMax{}
with median ARI \cvMedianARI{} against the adopted partition, better than the superseded default
(ARI \oldMedianARI{}, \oldClusters{} clusters, \oldNoisePct\% unassigned) but short of a
subfamily claim. \emph{Reproducibility:} clustering is deterministic within an environment (ARI
$1.000$ on re-run), but UMAP runs with \texttt{n\_jobs}$=-1$, so counts can shift with core
count; every figure uses the stored labels and embedding. \emph{Consensus cross-check:}
projecting each sampled copy onto every per-cluster consensus with
\texttt{mafft --addfragments --keeplength}, only \xcAccuracy\% are nearest their own cluster's
consensus (chance \xcChance\%) and the median copy is marginally closer to another's
(\xcDivOwn{} vs \xcDivOther{}), so we do not call the clusters subfamilies. What they capture is
degradation state: ORF integrity runs from \intactRich\% intact in the most intact cluster to
\intactPoor\% in the least, and expression concentrates there, which is the axis the workflow
needs.

The core workflow together with the phylogenetics, gRNA, transduction and
antisense modules is the source of every number reported above; the further analyses of
Section~\ref{sec:exploratory} extend the same resolved family into regulatory, evolutionary and
structural context.

\section{Design and implementation}
\label{sec:design}

TERRA is organised in three layers (Figure~\ref{fig:arch}). A Tauri/React desktop application
gives point-and-click access and performs no analysis itself; it launches a bundled Python
sidecar over newline-delimited JSON, so jobs stream progress and can be cancelled. The backend
is a driver (\texttt{query.py}) orchestrating the core stages, the standalone
\texttt{run\_*.py} modules and an SSH client: locally it fetches sequence from the UCSC
API~\citep{kent2002ucsc} when no FASTA is present, and on HPC it stages code, submits a batch
job and retrieves results. Every external resource (RMSK/Dfam,
UCSC~\citep{kent2002ucsc}, JASPAR~\citep{castromondragon2022jaspar},
mygene.info~\citep{xin2016mygene}) sits behind a thin adapter, so adding an assembly or source
is a localised change.

The pipeline is a directed sequence of stages (Table~\ref{tab:stages}). Each stage reads the
previous stage's tabular output and writes a checkpoint, so a run can be resumed or a single
stage re-executed. Figure~\ref{fig:flowchart} gives the end-to-end view.


The dominant cost in whole-family primer design is
reading the reference genome, since a naive implementation re-reads the multi-gigabyte FASTA
for every candidate search. TERRA parses the reference once into an in-memory genome cache,
reused for sequence extraction and all primer searches, and persists it so later runs skip
parsing. A small number of batch operations (GC content, length tabulation, $k$-mer matrix
construction) dominate the per-sequence arithmetic and have Cython~\citep{behnel2011cython}
implementations with a pure-Python fallback.

The driver can run the copy-resolved modules as a sequential
subprocess loop, or delegate them to a Nextflow DSL2 layer~\citep{ditommaso2017} that runs the
same unmodified Python modules as parallel, individually-resourced, resumable tasks
(Figure~\ref{fig:dag}). No scientific code is duplicated between the two.


% ===========================================================================
\section{Additional analysis modules}
\label{sec:exploratory}

TERRA ships further modules that extend the resolved family into its regulatory, evolutionary
and structural context (Figure~\ref{fig:stage11overview}); each runs as an independent task on
the same per-copy tables and emits its own figures, values and provenance manifest.
\emph{Regulatory:} CTCF-site/TAD-boundary proximity and an epigenetic/regulatory-track overlay
separate silenced from active insertions, complementing the expression stratification of
Section~\ref{sec:resolve}. \emph{Evolutionary:} orthologous-insertion calling and multi-assembly
liftover (\texttt{hg38}, T2T-CHM13, \texttt{mm10}, \texttt{mm39}) place each copy across
assemblies and species, and automatic subfamily resolution collapses the per-cluster consensuses
into a dendrogram. \emph{Structural:} an LTR structural classifier labels each element
(full-length / solo-LTR / internal-only / truncated) and ColabFold folds the consensus ORFs.
Each runs on any family the core workflow accepts and writes the same output as the core stages.


% ===========================================================================
\section{Generality across TE classes}
\label{sec:supp}

To show the workflow is not LINE-specific, we ran the identical pipeline on a SINE and an LTR
family in \texttt{\protoAssembly} with the same six-stage SQuIRE dataset. \emph{B1\_Mus2}
(SINE) has \bOneMTLoci{} loci (\bOneMTExprLoci{} expressed, median divergence
\bOneMTMedianDiv\%) and \emph{IAPLTR1\_Mm} (LTR) has \iapLoci{} loci (\iapExprLoci{} expressed,
\iapMedianDiv\%). The divergence medians order the families by age, \emph{IAPLTR1\_Mm} youngest,
\emph{\protoFamily} intermediate (\protoMedianDiv\%) and \emph{B1\_Mus2} oldest, consistent with
their retrotransposition histories~\citep{bourque2018}. Full per-family reports (clustering,
expression, alignment, consensus, motif, primer and guide outputs) are in the Supplementary
Files.

% ===========================================================================
\section{Discussion}

TERRA fills a specific gap: the accessible, reproducible, cluster-ready path from a named TE
family to reagents that can be ordered. It complements rather than replaces discovery tools
such as RepeatModeler2~\citep{flynn2020}, expression quantifiers such as
SQuIRE~\citep{yang2019squire} and TEtranscripts~\citep{jin2015tetranscripts}, and general
platforms such as Galaxy~\citep{galaxy2022} and nf-core~\citep{ewels2020}.

Its contributions are threefold: conceptually, resolving a family into sub-groups and
intersecting them with locus-level expression lets a signal carried by a minority of copies
survive, and motivates the expressed-versus-common consensus distinction
(Section~\ref{sec:resolve}); practically, treating primers and allele-aware CRISPRa/i guides as
the deliverable and reporting the guide trade-off as a frontier rather than one score
(Section~\ref{sec:crispr}); and in engineering, the in-memory genome cache, compiled hot paths
and automatic LSF/Slurm support make whole-family analysis practical from laptop to cluster.

Limitations bound the system. Expression quality and its multi-mapping policy are inherited from
the upstream quantifier, and family resolution depends on external services (RMSK/Dfam, UCSC,
JASPAR, mygene.info), so a run is only as reproducible as the database releases it draws on;
supported assemblies are currently the common human and mouse builds. The $k$-mer clustering,
though fast and alignment-free, does not yield stable sub-family boundaries at the default
minimum cluster size (agreement near chance, median ARI \cvMedianARI{}), so we report the
partition as a device for separating expressed from inert copies, not a subfamily assignment;
and the guide frontier reports sequence-level predictions that require experimental
confirmation.

Natural extensions include per-locus epigenomic context (WGBS, ATAC-seq) to separate silenced
from derepressed copies, correlating per-cluster activity with nearby gene expression, somatic-
insertion calling from user BAMs, broader assembly support, and negative-binomial modelling of
count overdispersion~\citep{love2014deseq2,robinson2010edger}.

% ===========================================================================
\section*{Code availability}
TERRA is open source at \url{https://github.com/anmol-dash/te-scraper-and-analysis}; the
pipeline scripts there generate every figure and measured value cited here.


\clearpage
\appendix
\section{Supplementary figures and tables}
\label{sec:appendix}

The figures and table below support the results and validation reported in the main text. They are placed here to keep the main narrative on the prototype walk-through and the reagent deliverables; each is referenced from the section it belongs to.

\begin{table}[h]
  \centering
  \small
  \begin{tabular}{@{}lll@{}}
    \toprule
    \textbf{Stage} & \textbf{Function} & \textbf{Key methods / tools} \\
    \midrule
    Prep        & Loci, sequences, divergence & RMSK/Dfam~\citep{hubley2016}; FASTA or UCSC~\citep{kent2002ucsc}; Kimura milliDiv \\
    Clustering  & Sub-group the copies        & $k$-mers $\to$ SVD $\to$ UMAP~\citep{mcinnes2018umap} $\to$ HDBSCAN~\citep{mcinnes2017hdbscan} \\
    Alignment   & Consensus per group         & MAFFT~\citep{katoh2002mafft} + CIAlign~\citep{tumescheit2022cialign} \\
    Motif       & Enriched TFBS               & JASPAR~\citep{castromondragon2022jaspar} $\cap$ loci (BEDTools~\citep{quinlan2010bedtools}) + Fisher \\
    GO          & Gene/ontology context       & mygene.info~\citep{xin2016mygene} \\
    Expression  & Per-group profiles + DE     & per-locus counts (SQuIRE~\citep{yang2019squire}); Wilcoxon + BH \\
    Primers     & Family/locus primers        & $k$-mer candidates + genome-wide specificity \\
    gRNA        & Allele-aware CRISPRa/i      & PAM-aware candidates + seed-anchored off-target scoring \\
    \bottomrule
  \end{tabular}
  \caption{The TERRA core stages. Each stage consumes a tabular checkpoint and emits one,
    allowing partial re-runs. The gRNA and primer stages produce the deliverables of
    Section~\ref{sec:crispr}.}
  \label{tab:stages}
\end{table}

\begin{figure}[htbp]
  \centering
  \figorbox[0.70\textwidth]{fig_locus_filtering}
  \caption{Resolving copies instead of lumping them. Left: a family-level summary averages
    over meaningful, low-function and inert loci, so signal carried by a minority of copies is
    diluted. Right: TERRA partitions copies by sequence and intersects them with locus-level
    expression, separating the populations and distinguishing the expressed consensus from the
    common consensus. This panel is a schematic of the concept; the measured partition for
    \emph{\protoFamily} is in Figure~\ref{fig:proto_results}.}
  \label{fig:filtering}
\end{figure}

\begin{figure}[htbp]
  \centering
  \figorbox[0.60\textwidth]{fig_phylo_master}
  \caption{Master-element nomination for \emph{\protoFamily}. Copies are scored on ORF
    integrity against divergence from the consensus; \protoIntact{} copies are intact and the
    top-ranked candidate is \protoMaster{}.}
  \label{fig:phylo_master}
\end{figure}

\begin{figure}[htbp]
  \centering
  \figorbox[0.70\textwidth]{fig_cluster_search}
  \caption{Clustering parameter search for \emph{\protoFamily}. Each point is one of
    \csNConfigs{} configurations of $k$, \texttt{n\_neighbors}, \texttt{min\_cluster\_size} and
    \texttt{min\_samples}, plotted as the number of clusters recovered against the fraction of
    loci left unassigned. The pipeline's default minimum cluster size
    ($\lfloor N/\protoDivisor \rfloor = \oldMinClusterSize$) admits only \oldClusters{} clusters
    and leaves \oldNoisePct\% of the family unassigned; the selected configuration (star)
    recovers \protoClusters{} clusters with \protoNoisePctMeasured\% unassigned.}
  \label{fig:cluster_search}
\end{figure}

\begin{figure}[htbp]
  \centering
  \figorbox[0.70\textwidth]{fig_de_l1mdt}
  \caption{Differential expression between \emph{\protoFamily} clusters. Heatmap of
    $-\log_{10}(q)$ from pairwise Wilcoxon rank-sum tests with Benjamini--Hochberg correction;
    \protoDESig{} of \protoDETotal{} cluster-pair by stage comparisons reach $q<0.05$. Cluster-resolved testing recovers stage-specific
    differences that a family-level summary averages away (Section~\ref{sec:resolve}).}
  \label{fig:de}
\end{figure}

\begin{figure}[htbp]
  \centering
  \figorbox[0.70\textwidth]{fig_phylo_tree}
  \caption{Subfamily phylogeny for \emph{\protoFamily}. Neighbour-joining tree relating
    cluster consensuses and representative copies, built from the MAFFT alignment.}
  \label{fig:phylo_tree}
\end{figure}

\begin{figure}[htbp]
  \centering
  \figorbox[0.70\textwidth]{fig_crispr_deliverable}
  \caption{Allele-aware CRISPRa/i guide design for a repetitive family. Candidate guides are
    enumerated PAM-aware from the consensus and per-copy sequence, scored against the full
    copy landscape with seed-anchored, mismatch-tolerant matching, and reported as a
    coverage-versus-off-target trade-off. This panel is a schematic of the logic; the measured
    frontier for \emph{\protoFamily} is Figure~\ref{fig:grna_pareto}.}
  \label{fig:crispr}
\end{figure}

\begin{figure}[htbp]
  \centering
  \figorbox[0.70\textwidth]{fig_transduction_groups}
  \caption{3' transduction lineages for \emph{\protoFamily}: \protoLineages{} candidate
    lineages grouped by shared downstream-tail $k$-mers absent from the family consensus;
    \protoPolyA{} copies carry an AATAAA poly-A signal.}
  \label{fig:transduction}
\end{figure}

\begin{figure}[htbp]
  \centering
  \figorbox[0.70\textwidth]{fig_antisense_motifs}
  \caption{Antisense/bidirectional promoter content for \emph{\protoFamily}: core promoter
    elements scanned on both strands of each 5' region; \protoBidir{} of \protoLoci{} copies
    are flagged bidirectional.}
  \label{fig:antisense}
\end{figure}

\begin{figure}[htbp]
  \centering
  \figorbox[0.70\textwidth]{fig_cluster_crosscheck}
  \caption{Cross-check of the clusters against the consensus alignments for
    \emph{\protoFamily}. \textbf{Left:} fraction of each cluster's copies whose nearest consensus
    is each cluster's, after projecting every copy onto every consensus with
    \texttt{mafft --addfragments --keeplength}. A subfamily assignment would put the mass on the
    diagonal; the observed \xcAccuracy\% (chance \xcChance\%) does not. \textbf{Right:} Kimura
    divergence of each copy to its own cluster's consensus versus to the nearest other
    consensus; the distributions overlap, with the median copy marginally closer to another
    cluster's consensus.}
  \label{fig:crosscheck}
\end{figure}

\begin{figure}[htbp]
  \centering
  \figorbox[0.70\textwidth]{fig_cluster_validation}
  \caption{Clustering and noise-class validation for \emph{\protoFamily}. \textbf{Left:}
    Adjusted Rand Index of each $k$ and \texttt{n\_neighbors} configuration against the
    configuration the automatic search accepted. Values near 1 would indicate the partition is
    independent of the parameter choice; the observed median of \cvMedianARI{} indicates it is
    not, which is why these sub-family boundaries are not interpreted as biological.
    \textbf{Right:} HDBSCAN noise-class fraction per configuration, with the accepted
    configuration highlighted. Noise-class loci are retained as a labelled group and excluded
    from consensus building.}
  \label{fig:cluster_validation}
\end{figure}

\begin{figure}[htbp]
  \centering
  \figorbox[0.70\textwidth]{fig_architecture}
  \caption{TERRA's three-layer architecture. A Tauri/React desktop shell drives a Python
    analysis sidecar over a streaming NDJSON protocol; the backend runs locally or stages and
    submits jobs to an LSF/Slurm cluster. Coordinate, sequence, motif and annotation data are
    pulled from RMSK/Dfam, UCSC, JASPAR and mygene.info through uniform adapters.}
  \label{fig:arch}
\end{figure}

\begin{figure}[htbp]
  \centering
  \figorbox[0.70\textwidth]{fig_pipeline_flowchart}
  \caption{TERRA pipeline flowchart. A family/assembly request flows through the core stages
    and into the copy-resolved modules, which run in parallel; outputs are gathered into a
    per-family folder and report. The run executes inside the Nextflow DSL2 and Singularity
    environment on an LSF or Slurm scheduler.}
  \label{fig:flowchart}
\end{figure}

\begin{figure}[htbp]
  \centering
  \figorbox[0.70\textwidth]{fig_nextflow_dag}
  \caption{The Nextflow DSL2 execution graph: the core engine feeds a scatter of independent
    copy-resolved module tasks, gathered into a single per-family output. Profiles select LSF
    or Slurm and Singularity or Docker.}
  \label{fig:dag}
\end{figure}

\begin{figure}[htbp]
  \centering
  \figorbox[0.70\textwidth]{fig_stage11_analyses}
  \caption{The copy-resolved analysis modules grouped by biological theme. Each card names the
    standalone \texttt{run\_*.py} script. The phylogenetics, gRNA, transduction and antisense
    modules are reported as results in Section~\ref{sec:results}; the remainder are the
    additional analysis modules of Section~\ref{sec:exploratory}.}
  \label{fig:stage11overview}
\end{figure}

\clearpage
\bibliographystyle{unsrtnat}
\bibliography{references}

\end{document}
